% Part: first-order-logic
% Chapter: proof-systems
% Section: introduction

\documentclass[../../../include/open-logic-section]{subfiles}

\begin{document}

\iftag{FOL}
      {\olfileid[fr]{fol}{prf}{int}}
      {\olfileid[fr]{pl}{prf}{int}}

\olsection{Introduction}

Une logique comporte généralement une sémantique et un système de
!!{derivation}. La sémantique porte sur des notions telles que la vérité,
la satisfaisabilité, la validité et la conséquence logique. Les systèmes
de !!{derivation} ont pour objet de fournir une méthode purement
syntaxique pour établir la conséquence logique et la validité. Ils sont
purement syntaxiques au sens où une !!{derivation} dans un tel système
est un objet syntaxique fini, généralement une suite (ou un autre
agencement fini) d'!!{sentence}s ou de !!{formula}s. Un bon système de
!!{derivation} permet de vérifier mécaniquement si une suite ou un
agencement donné d'!!{sentence}s ou de !!{formula}s est « correct ».

Les systèmes de !!{derivation} les plus simples, et historiquement les
premiers, pour la logique du premier ordre étaient \emph{axiomatiques}.
Une suite de !!{formula}s constitue une !!{derivation} dans un tel
système si chacune de ses !!{formula}s appartient à un ensemble fixé
d'« axiomes » ou se déduit de !!{formula}s antérieures de la suite par
l'une d'un nombre fixé de « règles d'inférence ». On peut vérifier
mécaniquement si une !!{formula} est un axiome et si elle se déduit
correctement d'autres !!{formula}s par l'une de ces règles. Les systèmes
axiomatiques de !!{derivation} sont faciles à décrire et se prêtent bien
à l'étude métathéorique. Leurs !!{derivation}s sont toutefois difficiles
à lire, à comprendre et à construire.

D'autres systèmes ont été élaborés pour faciliter la construction des
!!{derivation}s ou leur compréhension une fois achevées. C'est le cas de
la déduction naturelle, des arbres de vérité, également appelés preuves
par tableaux, et du calcul des séquents. Certains systèmes de
!!{derivation} sont spécialement conçus pour l'automatisation : la méthode
de résolution, par exemple, est facile à programmer, mais ses
!!{derivation}s sont pratiquement impossibles à comprendre. La plupart
de ces autres systèmes représentent les !!{derivation}s par des arbres
de !!{formula}s plutôt que par des suites. On voit ainsi plus facilement
de quelles autres parties dépend chaque partie d'une !!{derivation}.

Pour une logique donnée, par exemple celle du premier ordre, les
différents systèmes de !!{derivation} précisent donc de différentes
manières ce que signifie être un \emph{théorème} pour un
!!{sentence}, ou être !!{derivable} à partir d'autres !!{sentence}s.
Quelle que soit la présentation retenue (!!{derivation}s axiomatiques,
déductions naturelles, !!{derivation}s en calcul des séquents, arbres de
vérité ou réfutations par résolution), nous voulons que ces relations
correspondent aux notions sémantiques de validité et de conséquence
logique. Écrivons $\Proves !A$ pour « $!A$ est un théorème » et
« $\Gamma \Proves !A$ » pour « $!A$ est !!{derivable} à partir
de~$\Gamma$ ». Quelle que soit la définition de~$\Proves$, nous voulons
qu'elle corresponde à~$\Entails$, c'est-à-dire :
\begin{enumerate}
\item $\Proves !A$ si et seulement si $\Entails !A$ ;
\item $\Gamma \Proves !A$ si et seulement si $\Gamma \Entails !A$.
\end{enumerate}
Le sens « seulement si » de ces équivalences s'appelle la
\emph{correction}. Un système de !!{derivation} est correct si la
!!{derivability} garantit la conséquence logique (ou la validité).
Tout système de !!{derivation} digne de ce nom doit être correct ; un
système incorrect ne sert à rien. En effet, une !!{derivation} a
précisément pour fonction de donner une garantie syntaxique de validité
ou de conséquence logique. Nous démontrerons la correction des systèmes
de !!{derivation} que nous présenterons.

Le sens réciproque, « si », est lui aussi important : il s'appelle la
\emph{complétude}. Un système de !!{derivation} complet est assez
puissant pour établir que $!A$ est un théorème dès que $!A$ est valide,
et que $\Gamma \Proves !A$ dès que $\Gamma \Entails !A$.
La complétude est plus difficile à établir, et certaines logiques
n'admettent aucun système de !!{derivation} complet. La logique du
premier ordre en admet. Kurt G\"odel fut le premier à démontrer la
complétude d'un système de !!{derivation} pour la logique du premier
ordre, dans sa thèse de 1929.

Une autre notion liée aux systèmes de !!{derivation} est celle de
\emph{cohérence}. Un ensemble d'!!{sentence}s est dit incohérent si
l'on peut en !!{derive} n'importe quel !!{sentence}, et cohérent dans
le cas contraire. L'incohérence est le pendant syntaxique de
l'insatisfaisabilité : comme les ensembles insatisfaisables, les
ensembles incohérents d'!!{sentence}s ne constituent pas de bonnes
théories ; ils présentent un défaut fondamental. Les ensembles cohérents
d'!!{sentence}s ne sont pas nécessairement vrais ni utiles, mais ils
satisfont au moins cette condition minimale d'utilité logique. La
définition précise de la cohérence d'un ensemble d'!!{sentence}s peut
varier selon le système de !!{derivation}. Comme pour~$\Proves$, nous
voulons cependant qu'elle coïncide avec son pendant sémantique, la
satisfaisabilité : $\Gamma$ doit être cohérent si et seulement s'il est
satisfaisable. Ici, le sens « seulement si » correspond à la complétude
(la cohérence garantit la satisfaisabilité), et le sens « si » à la
correction (la satisfaisabilité garantit la cohérence). En logique
classique du premier ordre, les deux versions de la correction et de
la complétude sont en fait équivalentes.

\end{document}
